% Week 7 section of the synthesis exercise sheet, included by Exercises.tex.
\section{Week 7 --- Martingale Pricing, Hedging, and Market Structure}

Throughout, \(\Qbb\) denotes a martingale measure for the stated numeraire.
Every time a local martingale is used inside a conditional expectation,
state the integrability or admissibility assumption that makes it a true
martingale.

\subsection*{Part I --- Discounted Gains and Conditional Valuation}

\begin{exercise}[Medium --- Discounted gains when one asset pays dividends]
Let
\[
B_t=\exp\left(\int_0^t r_u\dd u\right),
\qquad D_t=B_t^{-1}.
\]
Two risky assets are traded. Asset \(S^1\) pays no dividends, while one share
of \(S^2\) pays the continuous cashflow \(q_tS_t^2\dd t\). A predictable
self-financing strategy has holdings
\((\phi_t^0,\phi_t^1,\phi_t^2)\) and wealth
\[
X_t=\phi_t^0B_t+\phi_t^1S_t^1+\phi_t^2S_t^2.
\]
\begin{exparts}
  \item Write the correct self-financing gain equation for \(X\).
  \item Prove the discounted gain identity
  \[
  \dd(D_tX_t)
  =\phi_t^1\dd(D_tS_t^1)
  +\phi_t^2\left(\dd(D_tS_t^2)+D_tq_tS_t^2\dd t\right).
  \]
  \item Suppose that, under \(\Qbb\), the discounted gain of each risky
  asset is a local martingale. Explain why \(DX\) is then a local martingale.
  Which additional issue must be settled before writing
  \[
  D_tX_t=\E^{\Qbb}[D_TX_T\mid\F_t]?
  \]
  \item A student claims that \(DS^2\) itself must be a martingale. Explain
  precisely why this is generally false when \(q\neq0\).
\end{exparts}
\end{exercise}

\begin{exercise}[Medium --- A payoff determined at an intermediate date]
Assume \(r\) is constant and \(B_t=e^{rt}\). Fix \(0<u<T\), and let
\(Y\in L^1(\Qbb)\) be \(\F_u\)-measurable. At time \(T\), a contract pays
\[
H=e^{r(T-u)}Y.
\]
\begin{exparts}
  \item For \(t\leq u\), derive the price process in terms of
  \(\E^{\Qbb}[Y\mid\F_t]\).
  \item For \(u\leq t\leq T\), show that the price is
  \(V_t=e^{r(t-u)}Y\).
  \item Check that the two expressions agree at \(t=u\), and interpret the
  contract as an operation performed at time \(u\).
  \item Prove directly that \(e^{-rt}V_t\) is a martingale on \([0,T]\).
  Why is the tower property the key step before \(u\)?
\end{exparts}
\end{exercise}

\subsection*{Part II --- New Martingale Representations and Hedges}

\begin{exercise}[Medium --- An exponential terminal claim]
Let \(W\) be a \(\Qbb\)-Brownian motion, let \(\theta\in\R\), and assume
zero interest for simplicity. Consider
\[
H=\exp\left(\theta W_T-\frac12\theta^2T\right).
\]
\begin{exparts}
  \item Compute \(M_t=\E^{\Qbb}[H\mid\F_t]\) using independent Brownian
  increments.
  \item Apply It\^o's formula to identify the process \(\psi_t\) in
  \[
  M_t=M_0+\int_0^t\psi_s\dd W_s.
  \]
  \item Suppose the discounted traded stock satisfies
  \[
  \dd\bar S_t=\sigma_t\bar S_t\dd W_t,
  \qquad \sigma_t\bar S_t\neq0.
  \]
  Find the stock holding that replicates \(H\).
  \item What fails on a set where \(\sigma_t\bar S_t=0\)? Distinguish the
  existence of the Brownian representation from its implementation by
  traded assets.
\end{exparts}
\end{exercise}

\begin{exercise}[Medium --- The Brownian area as a hedgeable claim]
Let \(W\) be a \(\Qbb\)-Brownian motion, \(r=0\), and
\[
H=\int_0^T W_s\dd s.
\]
\begin{exparts}
  \item Use stochastic Fubini or deterministic integration by parts to prove
  \[
  H=\int_0^T(T-u)\dd W_u.
  \]
  \item Compute \(M_t=\E^{\Qbb}[H\mid\F_t]\) in terms of
  \(\int_0^tW_s\dd s\) and \(W_t\).
  \item Apply It\^o's formula to your answer and recover the representation
  coefficient from part (a).
  \item Compute \(\Var^{\Qbb}(H)\) in two different ways.
  \item If \(\dd\bar S_t=\gamma_t\dd W_t\) with \(\gamma_t\neq0\), find the
  replicating stock holding. Why does the hedge decrease as maturity
  approaches?
\end{exparts}
\end{exercise}

\subsection*{Part III --- Spanning, Redundancy, and Claim-Specific Completeness}

\begin{exercise}[Medium --- Complete risk spanning with redundant assets]
Three discounted risky assets are driven by two \(\Qbb\)-Brownian factors:
\[
\dd\bar S_t=\Gamma\dd W_t,
\qquad
\Gamma=
\begin{pmatrix}
1&0\\
0&1\\
1&1
\end{pmatrix}.
\]
A discounted claim has MRT exposure
\(\psi_t=(a_t,b_t)^\top\).
\begin{exparts}
  \item Show that every \(\psi_t\in\R^2\) can be hedged.
  \item Find the full family of holdings
  \(\phi_t=(\phi_t^1,\phi_t^2,\phi_t^3)^\top\) satisfying
  \(\Gamma^\top\phi_t=\psi_t\).
  \item Among this family, find the hedge with the smallest Euclidean norm.
  \item Identify a nonzero zero-exposure portfolio. Explain why nonuniqueness
  of holdings here does not imply market incompleteness or price ambiguity.
\end{exparts}
\end{exercise}

\begin{exercise}[Medium to challenging --- Which claims remain replicable?]
The filtration is generated by two independent \(\Qbb\)-Brownian motions
\((W^1,W^2)\). The only discounted risky asset satisfies
\[
\dd\bar S_t=\sigma_t\bar S_t\dd W_t^1,
\qquad \sigma_t\bar S_t\neq0.
\]
Let deterministic \(f,g\in L^2(0,T)\), and consider
\[
H_{f,g}
=c+\int_0^Tf_t\dd W_t^1+\int_0^Tg_t\dd W_t^2.
\]
\begin{exparts}
  \item Find the martingale
  \(M_t=\E^{\Qbb}[H_{f,g}\mid\F_t]\) and its two-dimensional MRT exposure.
  \item Prove that \(H_{f,g}\) is replicable using the stock and cash if and
  only if \(g=0\) almost everywhere.
  \item When \(g=0\), find the stock and cash holdings that replicate the
  claim.
  \item Explain why the market can be incomplete even though every claim in
  the subfamily \(\{H_{f,0}\}\) has a unique replication price.
\end{exparts}
\end{exercise}


\subsection*{Part IV --- Forward Pricing and Unspanned Risk}

\begin{exercise}[Medium --- A forward with time-dependent deterministic rates]
Assume under \(\Qbb\)
\[
\frac{\dd S_t}{S_t}=r(t)\dd t+\sigma(t)\dd W_t^{\Qbb},
\]
where \(r,\sigma\) are deterministic and square-integrable on \([0,T]\).
Let \(B_t=\exp(\int_0^tr(u)\dd u)\).

\begin{exparts}
  \item Apply It\^o's formula to \(\log S_u\), integrate from \(t\) to
  \(T\), and exponentiate to obtain \(S_T\). Then give the conditional
  mean and variance of \(\log(S_T/S_t)\) given \(\F_t\). Why is its
  stochastic-integral term independent of \(\F_t\)?
  \item Prove directly that \(S_t/B_t\) is a true \(\Qbb\)-martingale.
  State the integrability fact you use.
  \item Price the forward-style payoff \(H=S_T-K\) by conditional
  expectation and show
  \[
  V_t=S_t-K\exp\!\left(-\int_t^T r(u)\dd u\right).
  \]
  \item Determine the replicating stock holding and explain why the
  volatility function disappears from the price but not from the
  distribution of the payoff.
\end{exparts}
\end{exercise}

\begin{exercise}[Medium to challenging --- An unspanned Brownian risk]
Let \(W^1,W^2\) be independent Brownian motions under \(\Pb\). The only
risky traded asset satisfies
\[
\frac{\dd S_t}{S_t}=\mu\dd t+\sigma\dd W_t^1,
\qquad \sigma>0,
\]
and the short rate is \(r\). Set \(\lambda=(\mu-r)/\sigma\). For each
\(\eta\in\R\), define
\[
Z_T^\eta=
\exp\!\left(
-\lambda W_T^1-\eta W_T^2
-\frac12(\lambda^2+\eta^2)T
\right).
\]
\begin{exparts}
  \item Verify that \(Z^\eta\) is a density process and identify the two
  Brownian motions under the associated measure \(\Qbb^\eta\).
  \item Show that \(e^{-rt}S_t\) is a martingale under every
  \(\Qbb^\eta\). Why does the traded stock impose no restriction on
  \(\eta\)?
  \item Consider the nontraded claim \(H=(W_T^2)^2\). Compute
  \(\E^{\Qbb^\eta}[H]\) and the corresponding discounted expectation price.
  \item Explain why the family of different prices does not contradict
  no-arbitrage. What missing traded risk would be needed to make the claim
  replicable?
\end{exparts}
\end{exercise}

\subsection*{Part V --- Path Dependence}

\begin{exercise}[Medium --- A weighted and floating-strike Asian family]
Under \(\Qbb\), assume
\[
\frac{\dd S_t}{S_t}=(r-q)\dd t+\sigma\dd W_t,
\]
with constant \(r,q,\sigma\). Let \(w:[0,T]\to[0,\infty)\) be deterministic
and define
\[
A_t=\int_0^t w(u)S_u\dd u.
\]
Consider the payoff
\[
H=(\alpha S_T+\beta A_T-K)^+.
\]
\begin{exparts}
  \item Show that \((S_t,A_t)\) is a Markov state and write its dynamics.
  \item Assuming \(V_t=v(t,S_t,A_t)\), derive the pricing PDE and terminal
  condition.
  \item Explain why there are no \(v_{aa}\) or \(v_{sa}\) terms.
  \item Identify the stock holding in the replicating hedge.
  \item Recover as special cases a weighted fixed-strike Asian call and a
  floating-strike Asian call. State the corresponding choices of
  \(\alpha,\beta,K\).
\end{exparts}
\end{exercise}

\subsection*{Part VI --- A Finite-Market Test of the FTAP}

\begin{exercise}[Challenging --- EMM interval, replication, and superhedging]
Consider a one-period market with zero interest and three states
\(\omega_1,\omega_2,\omega_3\), each having positive physical probability.
The bond satisfies \(B_0=B_1=1\), while
\[
S_0=2,
\qquad
S_1(\omega_1)=1,\quad
S_1(\omega_2)=2,\quad
S_1(\omega_3)=4.
\]
\begin{exparts}
  \item Find every equivalent martingale measure
  \(\Qbb^c\). Express its three state probabilities in terms of one
  parameter \(c\) and determine the admissible interval for \(c\).
  \item Use the FTAP to explain why the market is arbitrage-free but
  incomplete.
  \item Show that a payoff \(h=(h_1,h_2,h_3)\) is replicable if and only if
  \[
  h_3-3h_2+2h_1=0.
  \]
  \item Let \(H=(0,1,0)\), the claim paying only in the middle state.
  Compute its expectation under every \(\Qbb^c\) and determine the full
  no-arbitrage price interval when \(H\) is added as a traded claim.
  \item Find the cheapest stock--bond superhedge and the most expensive
  stock--bond subhedge for \(H\). Relate their costs to the endpoints of the
  price interval.
\end{exparts}
\end{exercise}
